\documentclass[10pt]{mypackage}

\usepackage[uprightgreeks,noDcommand]{kpfonts}
\renewcommand{\coloneq}{\coloneqq}
\renewcommand{\eqcolon}{\eqqcolon}
%\usepackage{dsfont}
%\renewcommand{\mathbb}{\mathds}
\usepackage{homework}
%\usepackage{notes}

%\usepackage[ backend=bibtex, style = alphabetic, sorting=ynt ]{biblatex}
%\addbibresource{  }

\usepackage{parskip}
\frenchspacing

\fancyhf{}
\fancyhead[R]{Avinash Iyer, Shoya Shimazaki, Liam Staker}
\fancyhead[L]{Functional Analysis: Homework 3}
\fancyfoot[C]{\thepage}

\setcounter{secnumdepth}{0}

\begin{document}
\RaggedRight
\section{Problem I}%

\section{Problem II}%

\section{Problem III}%

\section{Problem IV}%
Let $\left( X,\mathcal{T}_X \right)$ be a nonempty topological space. Consider $X^{\N}$ equipped with the product topology. For any $N\in \N$, we will also consider $X^{N+1}$ equipped with the product topology. For any $M\leq N$, we introduce the maps
\begin{align*}
  \Pi_{N,M}^X\colon X^{N+1} &\to X^{M+1}\\
  \left( x_0,\dots,x_N \right) &\mapsto \left( x_0,\dots,x_M \right)
\end{align*}
and
\begin{align*}
  \pi_{N,M}^{X}\colon X^{N+1} &\to X\\
  \left( x_0,\dots,x_N \right) &\to X_M.
\end{align*}
We will also consider their $N = \infty$ analogues; for any $M\in \N$, we let
\begin{align*}
  \Pi_{\infty,M}^X\colon X^{\N} &\to X^{M+1}\\
  \left( x_0,x_1,\dots \right) &\to \left( x_0,\dots,x_M \right)\\
  \pi_{\infty,M}^{X}\colon X^{\N} &\to X\\
  \left( x_0,x_1,\dots \right) &\to x_M.
\end{align*}
All these maps are continuous by the definition of the product topology. We will call a sequence $\left( \mu_N \right)_{N\geq 0}$ with $\mu_N\in \operatorname{Prob}\left( X^{N+1} \right)$ \textit{coherent} if for all $M,N\in \N$ with $M\leq N$, we have that $\left( \Pi^{X}_{N,M} \right)_{\ast}\mu_N = \mu_M$.

We will say that $\left( X,\mathcal{T}_X \right)$ satisfies the Daniell--Kolmogorov extension theorem (KET) if for all coherent sequences of probability measures $\left( \mu_N \right)_N$, there is a unique $\nu\in \operatorname{Prob}\left( X^{\N} \right)$ such that for all $N\geq 0$, one has $\left( \Pi^{X}_{\infty,N} \right)_{\ast}\nu = \mu_N$.
\begin{problem}[Problem IV.1]
  If $X$ and $Y$ are continuous spaces, and $\theta\colon X\to Y$ is a map between them, then for any $N\geq 0$ we let
  \begin{align*}
    \theta^{(N)}\colon X^{N+1} &\to Y^{N+1}\\
    \left( x_0,\dots,x_N \right) &\to \left( \theta\left( x_0 \right),\dots,\theta\left( x_N \right) \right),
  \end{align*}
  and analogously define the infinite analogue $\theta^{(\infty)}\colon X^{\N}\to Y^{\N}$. Suppose $\theta$ is a homeomorphism. Show that $X$ satisfies the KET if and only if $Y$ does.
\end{problem}
\begin{solution}
  Suppose $X$ satisfies the KET, and let $\left( \mu_N \right)_{N}$ be a coherent sequence of probability measures on $Y$. We will let $\Phi = \theta^{-1}$, which is a homeomorphism from $Y$ to $X$. Consider the sequence of pushforwards $\left( \Phi^{(N)}_{\ast}\mu_N \right)_{N}$. We claim that this sequence is coherent with respect to $X$. This amounts to showing that for any $M\leq N$, we have that $\left( \Pi^X_{N,M} \right)_{\ast}\left( \Phi^{(N)}_{\ast}\mu_N \right) = \Phi^{(M)}_{\ast}\mu_M$.

  Let $E = E_0\times \cdots\times E_M$ be a Borel cylinder set in $X^{M+1}$, and let $F_i = \theta\left( E_i \right)$. Then, we have that
  \begin{align*}
    \Phi^{(M)}_{\ast}\mu_M\left( E_0\times\cdots\times E_M\right) &= \mu_M\left( \theta^{(M)}\left( E_0\times\cdots\times E_M \right) \right)\\
                                                                  &= \mu_M\left( \theta\left( E_0 \right)\times\cdots\times \theta\left( E_M \right) \right)\\
                                                                  &= \mu_M\left( F_1\times\cdots\times F_M \right),
  \end{align*}
  while
  \begin{align*}
    \left( \Phi^X_{N,M} \right)_{\ast}\left( \Phi^{(N)}_{\ast}\mu_N \right)\left( E_0\times\cdots\times E_M \right) &= \mu_N\left( \theta^{(N)}\left( E_0\times\cdots\times E_M\times X^{N-M} \right) \right)\\
                                        &= \mu_N\left( \theta\left( E_0 \right)\times\cdots\times \theta\left( E_M \right)\times Y^{N-M} \right)\\
                                        &= \mu_N\left( F_0\times\cdots\times F_M\times Y^{N-M} \right)\\
                                        &= \mu_N\left( \left( \Pi^Y_{N,M} \right)^{-1}\left( F_0\times\cdots\times F_M \right) \right)\\
                                        &= \mu_M\left( F_0\times\cdots\times F_M \right),
  \end{align*}
  where the latter equality emerges from the fact that the $\left( \mu_N \right)_{N}$ are coherent on $Y$. Since $X$ satisfies the KET and $\left( \Phi_{\ast}\mu_N \right)_{N}$ is coherent on $X$, there is a unique probability measure $\nu$ on $X^{\N}$ such that $\left( \Pi^{X}_{\infty,M} \right)_{\ast}\nu = \Phi^{(M)}_{\ast}\mu_{M}$. We claim that $\theta_{\ast}^{(\infty)}\nu$ is the desired probability measure on $Y^{\N}$. This amounts to verifying that $\left( \Pi^Y_{\infty,M} \right)_{\ast}\theta_{\ast}^{(\infty)}\nu = \mu_M$ on Borel cylinder sets in $Y^{\N}$, as the cylinder sets generate the Borel $\sigma$-algebra on $Y^{\N}$ and any measure is uniquely determined by its value on the cylinder sets by the Carathéodory extension theorem.

  Let $F_0\times\cdots\times F_M\times Y^{\N}$ be a Borel cylinder set in $Y^{\N}$. Write $E_i = \Phi\left( F_i \right)$. Then, we have that
  \begin{align*}
    \left( \Pi^{Y}_{\infty,M} \right)_{\ast}\left( \theta_{\ast}^{(\infty)}\nu \right)\left( F_0\times\cdots\times F_M\times Y^{\N} \right) &= \left( \Pi^{Y}_{\infty,M}\circ \theta^{(\infty)} \right)_{\ast}\nu\left( F_0\times\cdots\times F_M\times Y^{\N} \right)\\
                                                     &= \left( \theta^{(M)}\circ \Pi_{\infty,M}^{X} \right)_{\ast}\nu\left( F_0\times\cdots\times F_M\times Y^{\N} \right)\\
                                                     &= \left( \theta^{(M)} \right)_{\ast}\left( \Pi^{X}_{\infty,M} \right)_{\ast}\nu\left( F_0\times\cdots\times F_M\times Y^{\N} \right)\\
                                                     &= \theta^{(N)}_{\ast}\Phi^{(N)}_{\ast}\mu_M\left( F_0\times\cdots\times F_M \right)\\
                                                     &= \left( \theta^{(N)}\circ \Phi^{(N)} \right)_{\ast}\mu_M\left( F_0\times\cdots\times F_M \right)\\
                                                     &= \mu_M\left( F_0\times\cdots\times F_M \right).
  \end{align*}
  The reverse direction follows from swapping all instances of $\theta$ with instances of $\Phi$ and all instances of $X$ with instances of $Y$.
\end{solution}
\begin{problem}[Problem IV.2]
  We say that $\left( X,\mathcal{T}_X \right)$ is well-behaved with respect to continuous functions if for all $N\geq 0$ and for all $\mu_1,\mu_2\in \operatorname{Prob}\left( X^{N+1} \right)$, we have that if
  \begin{align*}
    \int_{X^{N+1}}^{} f\:d\mu_1 &= \int_{X^{N+1}}^{} f\:d\mu_2
  \end{align*}
  for all $f\in C^b\left( X^{N+1},\R \right)$, then $\mu_1 = \mu_2$. Show that if $X$ is compact metrizable, then $X$ is well-behaved with respect to continuous functions.
\end{problem}
\begin{solution}
  Suppose $N\geq 0$ and $\mu_1,\mu_2\in \operatorname{Prob}\left( X^{N+1} \right)$ are arbitrary. Notice that if $X$ is compact metrizable, then so is $X^{N+1}$ (choose a compatible metric on $X$, then take the coordinatewise sum), and so any continuous function is bounded. In particular, $C\left( X^{N+1} \right) = C^b\left( X^{N+1} \right)$. If we define linear functionals
  \begin{align*}
    \varphi_1 &= \int_{X^{N+1}}^{} \cdot\:d\mu_1\\
    \varphi_2 &= \int_{X^{N+1}}^{} \cdot\:d\mu_2,
  \end{align*}
  then we observe that we have that $\varphi_1(f) = \varphi_2(f)$ for all $f\in C\left( X^{N+1} \right)$, so that $\varphi_1 = \varphi_2$. By the Riesz representation theorem, since the measures $\mu_i$ are in one-to-one correspondence with $\varphi_i$, it follows that $\mu_1 = \mu_2$.
\end{solution}
\begin{problem}[Problem IV.3]
  Show that $X = \R$ is well-behaved with respect to continuous functions.
\end{problem}
\begin{solution}
  
\end{solution}
\begin{problem}[Problem IV.4]
  Show that if $\theta\colon X\to Y$ is a homeomorphism, then $X$ is well-behaved with respect to continuous functions if and only if $Y$ is well-behaved with respect to continuous functions.
\end{problem}
\begin{solution}
  Suppose $X$ is well-behaved with respect to continuous functions. Let $\mu_1,\mu_2\in \operatorname{Prob}\left( Y \right)$ be probability measures, and suppose that for all $f\in C^b\left( Y^{N+1},\R \right)$, one has that
  \begin{align*}
    \int_{Y^{N+1}}^{} f\:d\mu_1 &= \int_{Y^{N+1}}^{} f\:d\mu_2.
  \end{align*}
  If we let $\theta^{(N+1)}\colon X^{N+1}\to Y^{N+1}$ be the corresponding amplified homeomorphism, then one observes by the way the pushforward measure is defined that
  \begin{align*}
    \int_{X^{N+1}}^{} f\circ \theta^{(N+1)}\:d\theta^{(N+1)}_{\ast}\mu_1 &= \int_{X^{N+1}}^{} f\circ \theta^{(N+1)}\:d\theta^{(N+1)}_{\ast}\mu_2.
  \end{align*}
  Since $\theta^{(N+1)}$ is a homeomorphism, the map $f\mapsto f\circ \theta^{(N+1)}$ is a bijection between $C^b\left( Y^{N+1} \right)$ and $C^b\left( X^{N+1} \right)$ with inverse map $g\mapsto g\circ \left( \theta^{-1} \right)^{(N+1)}$. Since $X$ is well-behaved with respect to continuous functions, one observes that we must have $\theta^{(N+1)}_{\ast}\mu_1 = \theta^{(N+1)}_{\ast}\mu_2$ as measures on $X$. Yet, since $\theta^{(N+1)}$ is a homeomorphism, it is necessarily the case that for any Borel $E\subseteq Y^{N+1}$, we have a unique Borel $F\subseteq X^{N+1}$ such that $\theta^{(N+1)}\left( F \right) = E$, whence $\mu_1 = \mu_2$. Therefore, $Y$ is well-behaved with respect to continuous functions.

  The reverse direction emerges from replacing all instances of $\theta$ with instances of $\theta^{-1}$ and vice versa.
\end{solution}
\begin{problem}[Problem IV.5]
  We now assume that $\left( X,\mathcal{T}_X \right)$ is compact metrizable. Define the collection of cylindrical continuous functions
  \begin{align*}
    \mathcal{A} &= \set{f\in C\left( X^{\N} \right) : \text{there exists $N\geq 0$ and $g\in C\left( X^{N+1} \right)$ such that $f = g\circ \Pi_{\infty,N}^{X}$}}.
  \end{align*}
  Show that $\mathcal{A}$ is a unital (self-adjoint) subalgebra of $C\left( X^{\N} \right)$.
\end{problem}
\begin{solution}\hfill
  \begin{itemize}
    \item Notice that if $f\in \mathcal{A}$, then $ \overline{f} = \overline{g}\circ \Pi_{\infty,N}^{X} $ so that $\mathcal{A}$ is self-adjoint.
    \item We observe that the function $1\colon X^{\N}\to \C$ can be realized by the first coordinate --- i.e., $1 = 1\circ \Pi_{\infty,0}^{X}$. Therefore, $\mathcal{A}$ is unital.
    \item If $f,g\in \mathcal{A}$, then we have $M$ such that $\hat{f}\in C\left( X^{M+1} \right)$ satisfies $f = \hat{f}\circ \Pi_{\infty,M}^{X}$ and $N$ such that $\hat{g}\in C\left( X^{N+1} \right)$ satisfies $g = \hat{g}\circ \Pi^{X}_{\infty,N}$. Without loss of generality, assume that $M\leq N$. Then, we observe that $\Pi_{N,M}^{X}\circ \Pi_{\infty,M}^X = \Pi_{\infty,M}^{X}$, so that $f = \left( \hat{f}\circ \Pi_{N,M}^{X} \right)\circ \Pi_{\infty,N}^{X}$.

      In particular, this means that the functions $\hat{g}$ and $\hat{f}\circ \Pi^{X}_{N,M}$ are defined on $X^{N+1}$, so that we may write
      \begin{align*}
        f+g &= \left( \hat{g} + \hat{f}\circ \Pi^{X}_{N,M}\right)\circ \Pi^{X}_{\infty,N}\\
        fg &= \left( \left( \hat{g} \right)\left( \hat{f}\circ \Pi^{X}_{N,M} \right) \right)\circ \Pi^{X}_{\infty,N}.
      \end{align*}
      Thus, $\mathcal{A}$ is a subalgebra of $C\left( X^{\N} \right)$.
  \end{itemize}
\end{solution}
\begin{problem}[Problem IV.6]
  Use Stone--Weierstrass to show that $ \mathcal{A} $ is uniformly dense in $C\left( X^{\N} \right)$.
\end{problem}
\begin{solution}
  It suffices to show that $\mathcal{A}$ separates points. If $\left( x_n \right)_{n}$ and $\left( y_n \right)_{n}$ are distinct elements of $X^{\N}$, then let $N$ be the smallest index for which $x_N\neq y_N$. Then, $\left( x_0,\dots,x_N \right)\neq \left( y_0,\dots,y_N \right)$, so since $X^{N+1}$ is compact metrizable, there is a continuous function $\hat{f}\colon X^{N+1}\to \C$ such that $\hat{f}\left( x_0,\dots, x_N\right)\neq \hat{f}\left( y_0,\dots,y_N \right)$. Thus, $f = \hat{f}\circ \Pi^{X}_{\infty,N}$ takes $\left( x_n \right)_{n}$ and $\left( y_n \right)_{n}$ to different values. Thus, $\mathcal{A}$ is a unital self-adjoint subalgebra of $C\left( X^{\N} \right)$ that separates points, so it is dense.
\end{solution}
\begin{problem}[Problem IV.7]
  Let $\left( \mu_N \right)_{N\geq 0}$ be a coherent sequence of probability measures adapted to $X$. Define the map $L\colon \mathcal{A}\to \R$ by taking
  \begin{align*}
    L\left( f \right) &= \int_{X^{N+1}}^{} g\:d\mu_N
  \end{align*}
  whenever one may write $f = g\circ \Pi_{\infty,N}^{X}$ with $N\geq 0$ and $g\in C\left( X^{N+1},\R \right)$. Show that this map is well-defined an $\R$-linear.
\end{problem}
\begin{solution}
  Let $M\leq N$. Suppose that $g_1\in C\left( X^{N+1},\R \right)$ and $g_2\in C\left( X^{M+1},\R \right)$ are such that $f = g_1\circ \Pi^{X}_{\infty,N}$ and $f = g_2\circ \Pi^{X}_{\infty,M}$. Then, by the compatibility condition that $\Pi^{X}_{N,M}\circ \Pi^{X}_{\infty,N} = \Pi^{X}_{\infty,M}$, we have that $f = \left( g_2\circ \Pi^{X}_{N,M} \right)\circ \Pi^{X}_{\infty,N}$, so since the projection maps $\Pi^{X}_{\infty,N}$ are surjective (hence right-invertible), we have that $g_1 = g_2\circ \Pi^{X}_{N,M}$. In particular, we observe that
  \begin{align*}
    L(f) &= \int_{X^{N+1}}^{} g_1\:d\mu_N\\
         &= \int_{X^{N+1}}^{} g_2\circ \Pi^{X}_{N,M}\:d\mu_N\\
         &= \int_{X^{M+1}}^{} g_2\:d\left( \Pi^{X}_{N,M} \right)_{\ast}\mu_N\\
         &= \int_{X^{M+1}}^{} g_2\:d\mu_M,
  \end{align*}
  where the second-to-final equality emerges from the definition of the pushforward measure and the final equality from the coherence of the sequence $\left( \mu_N \right)_N$. Thus, we obtain that $L(f)$ is well-defined. Since $L(f)$ is an integral over the compact metrizable space $X^{N+1}$, it follows that $L(f)$ is $\R$-linear.
\end{solution}
\begin{problem}[Problem IV.8]
  Show that there is a unique continuous linear form $ \widetilde{L}\colon C\left( X^{\N},\R \right)\to \R $ such that $ \widetilde{L}|_{\mathcal{A}} = L $.
\end{problem}
\begin{solution}
  We notice that for any $f\in \mathcal{A}$, we have that
  \begin{align*}
    \left\vert L(f) \right\vert &\leq \sup_{x\in X^{N+1}}\left\vert g(x) \right\vert\\
                                &= \sup_{x\in X^{\N}} \left\vert f(x) \right\vert\\
                                &= \norm{f},
  \end{align*}
  where the second-to-final equality is from the fact that $f$ is a cylindrical function. Thus, $L$ is a bounded linear functional on the dense $\ast$-subalgebra $\mathcal{A}\subseteq C\left( X^{\N},\R \right)$, so that it extends by uniform continuity to the closure $ \widetilde{L}\colon C\left( X^{\N},\R \right)\to \R $ such that $ \widetilde{L}|_{\mathcal{A}} = L $.
\end{solution}
\begin{problem}[Problem IV.9]
  Using the Riesz Representation Theorem for $ \widetilde{L} $, finish the proof that a compact metrizable space must satisfy the KET.
\end{problem}
\begin{solution}
  Since $1\in \mathcal{A}$ and $\norm{L}\leq 1$ with
  \begin{align*}
    L\left( 1 \right) &= \int_{X}^{} 1\:d\mu_0\\
                      &= 1
  \end{align*}
  it follows that $ \widetilde{L} $, being an extension of $L$ with the same norm satisfying $ \widetilde{L}(1) = L(1) = 1 $, is a norm-one positive linear functional. Since $X^{\N}$ is compact metrizable, it follows by the Riesz Representation Theorem that there is a unique probability measure $\mu$ such that
  \begin{align*}
    \widetilde{L}(f) &= \int_{X^{\N}}^{} f\:d\mu
  \end{align*}
  for all $f\in C\left( X^{\N} \right)$.
\end{solution}
\begin{problem}[Problem IV.10]
  We now suppose that $\left( X,\mathcal{T}_X \right)$ satisfies the KET. Let $Y\subseteq \mathcal{B}_X$ be a nonempty Borel subset. Equip $Y$ with the induced topology $ \mathcal{T}_Y = \iota^{\ast}\left( \mathcal{T}_X \right) $ where $\iota\colon Y\to X$ denotes the canonical injection.

  Prove that $\mathcal{T}_{Y^{\N}} = \left( \iota^{(\infty)} \right)^{\ast}\left( \mathcal{T}_{X^{\N}} \right)$, and deduce that $\mathcal{B}_{Y^{\N}} = \left( \iota^{(\infty)} \right)^{\ast}\left( \mathcal{B}_{X^{\N}} \right)$.
\end{problem}
\begin{solution}
  If $\iota\colon Y\to X$ is the canonical injection, then we observe that $Y^{\N}$ injects into $X^{\N}$ by the $\N$-fold amplification 
  \begin{align*}
    \iota^{(\infty)}\left( \left( y_n \right)_{n} \right) &= \left( \iota\left( y_n \right) \right)_{n}.
  \end{align*}
  We notice that
  \begin{align*}
    \left( \iota^{(\infty)} \right)^{\ast}\left( \mathcal{T}_{X^{\N}} \right) &= \set{U\in Y^{\N}: \text{there exists $V\in \mathcal{T}_{X^{\N}}$ such that $V\cap Y = U$}}\\
                                      &= \set{U\in Y^{\N} : \text{there exist $V_1,\dots,V_n\in \mathcal{T}_X$ such that $U = \bigcap_{j=1}^{n}\left( \pi_{\infty,i_j}^X \right)^{-1}\left( Y\cap V_j \right)$}}\\
                                      &= \set{U\in Y^{\N} : \text{there exist $U_1,\dots,U_n\in \mathcal{T}_Y$ such that $U = \bigcap_{j=1}^{n}\left( \pi_{\infty,i_j}^Y \right)^{-1}\left( U_{i_j} \right)$}}\\
                                      &= \mathcal{T}_{Y^{\N}}.
  \end{align*}
  Since the product topology on $Y^{\N}$ is induced by the inclusion of $Y^{\N}$ into $X^{\N}$, it follows that the $\sigma$-algebra produced by $\iota^{(\infty)}$ is given by the Borel $\sigma$-algebra generated by the topology on $Y^{\N}$.
\end{solution}
\begin{problem}[Problem IV.11]
  Let $\left( \mu_N^Y \right)_{N\geq 0}$ be a coherent sequence of probability measures adapted to $Y$. For all $N\geq 0$, define $\mu_N^X = \left( \iota^{(N)} \right)_{\ast}\mu_N^y$. Prove that $\left( \mu_N^X \right)_{N\geq 0}$ is a coherent sequence of probability measures adapted to $X$.
\end{problem}
\begin{solution}
  Suppose $M\leq N$. Observe that we have $\iota^{(M)}\circ \Pi^{Y}_{N,M} = \Pi^{X}_{N,M}\circ \iota^{(N)}$, so by the coherence of $\left( \mu_N^Y \right)_{N\geq 0}$, we have
  \begin{align*}
    \mu^X_{M} &= \left( \iota^{(M)} \right)_{\ast}\mu_{M}^Y\\
              &= \left( \iota^{(M)}\circ \Pi^Y_{N,M} \right)_{\ast} \mu_N^Y\\
              &= \left( \Pi_{N,M}^{X}\circ \iota^{(N)} \right)_{\ast}\mu_N^Y\\
              &= \left( \Pi_{N,M}^{X} \right)_{\ast}\mu_N^X.
  \end{align*}
\end{solution}
\begin{problem}[Problem IV.12]
  Prove the uniqueness part of the KET for $Y$.
\end{problem}
\begin{solution}
  Since $X$ satisfies the KET, there is a unique $\nu\in \operatorname{Prob}\left( X^{\N} \right)$ such that $\left( \Pi^X_{\infty,N} \right)_{\ast}\nu = \mu_N^X = \left( \iota^{(N)} \right)_{\ast}\mu_N^Y$.

  Suppose there are two probability measures $\nu_1$ and $\nu_2$ such that $\left( \Pi^Y_{\infty,N} \right)_{\ast}\nu_1 = \left( \Pi^Y_{\infty,N} \right)_{\ast}\nu_2 = \mu^Y_N$. Then, we observe that
  \begin{align*}
    \left( \Pi^{X}_{\infty,N}\right)_{\ast}\left( \left( \iota^{(\infty)} \right)_{\ast}\nu_{1} \right) &= \mu_{N}^X,
  \end{align*}
  so that $\left( \iota^{(\infty)} \right)_{\ast}\nu_1 = \nu$, and analogously for $\left( \iota^{(\infty)} \right)_{\ast}\nu_2$. Yet, this means that for all Borel subsets $B\subseteq X^{\N}$, we have that
  \begin{align*}
    \nu\left( B \right) &= \nu_1\left( B\cap Y^{\N} \right)\\
                        &= \nu_2\left( B\cap Y^{\N} \right),
  \end{align*}
  so since $\mathcal{B}_{Y^{\N}} = \left( \iota^{(\infty)} \right)^{\ast}\mathcal{B}_{X^{\N}}$, it follows that $\nu_1\left( E \right) = \nu_2\left( E \right)$ for all Borel subsets $E\in \mathcal{B}_{Y^{\N}}$. So, $\nu_1 = \nu_2$.
\end{solution}
\begin{problem}[Problem IV.13]
  By writing
  \begin{align*}
    Y^{\N} &= \bigcap_{N\geq 0} \left( \pi^{X}_{\infty,N} \right)^{-1}\left( Y \right),
  \end{align*}
  show that $Y^{\N}\in \mathcal{B}_{X^{\N}}$.
\end{problem}
\begin{solution}
  Since $\pi^X_{\infty,N}$ is continuous for each $N$, it is Borel, so $\left( \pi_{\infty,N}^X \right)^{-1}\left( Y \right)$ is Borel for each $N\in \N$, so since
  \begin{align*}
    Y^{\N} &= \bigcap_{N\geq 0} \left( \pi^X_{\infty,N} \right)^{-1}\left( Y \right),
  \end{align*}
  it follows that $Y^{\N}$ is the countable intersection of Borel sets, hence $Y^{\N}\in \mathcal{B}_{X^{\N}}$.
\end{solution}
\begin{problem}[Problem IV.14]
  Since $\left( \mu^X_{N} \right)_{N\geq 0}$ is a coherent sequence of probability measures adapted to $X$, and the latter satisfies the KET, we have a uniquely defined Borel probability measure $\nu^X$ on $X^{\N}$ such that for all $N\geq 0$, $\mu^X_N = \left( \pi^X_{\infty,N} \right)_{\ast}\nu^X$. Show that $\nu^X\left( X^{\N}\setminus Y^{\N} \right) = 0$.
\end{problem}
\begin{solution}
  We observe that for all $N\geq 0$, we have that
  \begin{align*}
    \left( \Pi^{\infty}_{X,N} \right)_{\ast}\nu^X\left( X^{\N}\setminus Y^{\N} \right) &= \mu^X_N\left( X^{N}\setminus Y^{N} \right)\\
                                                                                       &= \mu^Y_N\left( \left( X^N\setminus Y^N \right)\cap Y^N \right)\\
                                                                                        &= 0.
  \end{align*}
  Then, since
  \begin{align*}
    X^{\N}\setminus Y^{\N} &= \bigcap_{N\geq 0} \left( \Pi^{X}_{\infty,N} \right)^{-1}\left( X^{N}\setminus Y^{N} \right),
  \end{align*}
  with $\nu^X\left( X^{\N} \right) = 1$, it follows by continuity from above that $\nu^X\left( X^{\N}\setminus Y^{\N} \right) = 0$.
\end{solution}
\begin{problem}[Problem IV.15]
  Let $N\geq 0$ and let $f\in C^b\left( Y^{N+1},\R \right)$. We let $ \widetilde{f} $ be the extension of $f$ to $X^{N+1}$ by zero. We let $ \widetilde{f\circ \Pi^Y_{\infty,N}} $ the extension of $f\circ \Pi^Y_{\infty,N}$ from $Y^{\N}$ to $X^{\N}$ by zero. Show that $ \widetilde{f\circ \Pi^Y_{\infty,N}} = \widetilde{f}\circ \Pi^X_{\infty,N} $ $\nu^X$-a.e.
\end{problem}
\begin{solution}
  Upon bootstrapping, we observe that by the proof of the KET, we have that for all Borel $B\subseteq X^{N+1}$ that
  \begin{align*}
    \int_{\left( \Pi^X_{\infty,N} \right)^{-1}\left( B \right)}^{} \widetilde{f\circ \Pi^Y_{\infty,N}}\:d\nu^X &= \int_{\left( \Pi^X_{\infty,N} \right)^{-1}\left(B\right)\cap Y^{\N}}^{} f\circ \Pi^Y_{\infty}\:d\left( \nu^X|_{Y^{\N}} \right)\\
                  &= \int_{B\cap Y}^{} f\:d\left( \mu^X_N|_{Y^{N+1}} \right)\\
                  &= \int_{B\cap Y}^{} f\:d\mu^Y_N,
  \end{align*}
  while
  \begin{align*}
    \int_{\left( \Pi^X_{\infty,N} \right)^{-1}\left( B \right)}^{} \widetilde{f}\circ \Pi^{Y}_{\infty,N}\:d\nu^X &= \int_{B}^{} \widetilde{f}\:d\mu^X_N\\
                                 &= \int_{B\cap Y}^{} f\:d\mu^Y_N.
  \end{align*}
  Since this holds for every Borel subset of $X^{N+1}$, and the Borel subsets of $Y^{N+1}$ are intersections of Borel subsets of $X^{N+1}$ with $Y$, it follows that the functions $ \widetilde{f\circ \Pi^Y_{\infty,N}} $ and $ \widetilde{f}\circ \Pi^{X}_{\infty,N} $ are equal on all the Borel cylinder subsets of $X^{\N}$, so that they are equal $\nu^X$-a.e. by the Carathéodory extension theorem.
\end{solution}
\begin{problem}[Problem IV.16]
  Let $\nu^Y$ be the restriction to $Y^{\N}$ of the measure $\nu^X$. Show that $\nu^Y$ gives the solution to the existence problem regarding the KET for $Y$. Namely, show that for all $N\geq 0$, $\mu^Y_N = \left( \Pi^Y_{\infty,N} \right)_{\ast}\nu^Y$. Conclude that $\R$ satisfies the KET.
\end{problem}
\begin{solution}
  We observe that for all Borel $B\subseteq Y^{\N}$, we have that 
  \begin{align*}
    \nu^Y(B) &= \nu^X\left( \iota^{(\infty)}\left( B \right) \right).
  \end{align*}
  Let $B\subseteq Y^{N+1}$ be a Borel subset. Then,
  \begin{align*}
    \left( \Pi^Y_{\infty,N} \right)_{\ast}\nu^Y\left(B\right) &= \nu^Y\left( \left( \Pi^Y_{\infty,N} \right)^{-1}\left( B \right) \right)\\
                                                              &= \nu^X\left( \iota^{(\infty)}\left( \Pi^Y_{\infty,N} \right)^{-1}\left( B \right) \right)\\
                                                              &= \nu^X\left( \left( \Pi^{X}_{\infty,N} \right)^{-1}\left( \iota(B) \right) \right)\\
                                                              &= \left( \Pi^{X}_{\infty,N} \right)_{\ast}\nu^X\left( \iota(B) \right)\\
                                                              &= \mu^X_{N}\left( \iota(B) \right)\\
                                                              &= \mu^Y_N\left( B\cap Y \right)\\
                                                              &= \mu^Y_N\left( B \right).
  \end{align*}
  Since we may view $\R$ as a Borel subset of its one-point compactification, which is itself a compact metrizable space, we may thus find that $\R$ satisfies the KET.
\end{solution}
\end{document}
